% \newpage
\section{Towards an Information System}

This section presents the conceptual \gls{erd}, the key modeling assumptions and cardinalities, and how the provided sample is used to populate the database.

\subsection{Entity-relationship Diagram}

Figure~\ref{fig:erd} illustrates the \gls{erd} of the \gls{imdb}. The \gls{erd} models the main concepts required to represent entertainment titles, people involved in them, their classifications, their alternative regional or linguistic titles, and the episode hierarchy of series.

\begin{figure}[ht]
    \centering
    \includegraphics[width=1\linewidth]{fig/ERD.pdf}
    \caption{\gls{erd} for \gls{imdb}.}
    \label{fig:erd}
\end{figure}

% \textcolor{red}{normailze Principal.role\_names}

% \paragraph{Titles and classification.}
The central entity set is \texttt{Title}, representing an audiovisual work. It stores core descriptive attributes such as identifier, primary and original title, adult flag, years, and runtime. Each \texttt{Title} is classified by exactly one \texttt{ContentType} (e.g., movie, short, TV series). An \texttt{IS-A} specialization is only introduced for \texttt{Series} and \texttt{Episode} because episodic content requires additional structure that cannot be captured by classification alone: a \texttt{Series} can have episodes, while an \texttt{Episode} belongs to exactly one series and is positioned by \texttt{season\_number} and \texttt{episode\_number}. The specialization is partial, meaning a \texttt{Title} may be neither a series nor an episode; in that case, it is still represented as a \texttt{Title} and classified through \texttt{ContentType}. Genres are modeled using the weak entity set \texttt{Title\_Genre} between \texttt{Title} and \texttt{Genre}.

% \paragraph{People and credited participation.}
The entity set \texttt{Talent} represents people involved in titles (cast and crew). To model credits, we use the weak entity set \texttt{Principal} connecting \texttt{Title}, \texttt{Talent}, and \texttt{Category} (e.g., actor, director). This design is required because participation is not a simple many-to-many relationship. It carries additional attributes, namely \texttt{order}, \texttt{job}, and \texttt{role\_names}. In addition, professions of a person are represented through \texttt{Talent\_Role}, a weak entity set between \texttt{Talent} and \texttt{Role}, again preserving \texttt{order} where present.

% \paragraph{Alternative titles and localization.}
Alternative and localized titles are modeled through the weak entity set \texttt{TitleAKA}, which depends on \texttt{Title} via the identifying relationship \texttt{hasAKA}. TitleAKA is identified by its owning Title and a local discriminator \texttt{order}. It stores the alternative title text and related metadata. TitleAKA may optionally be linked to \texttt{Language} and \texttt{Region}, enabling language-specific and region-specific exploration. TitleAKA may also be classified by \texttt{TitleType} through the weak entity set \texttt{TitleAKA\_TitleType}.

% \paragraph{Series and episode hierarchy.}
Episodic content is represented through the relationship \texttt{hasEpisode} between \texttt{Series} and \texttt{Episode}, with \texttt{season\_number} and \texttt{episode\_number} supporting ordered navigation.

A more detailed overview of the modeling assumptions, cardinalities, and attribute domains used in the relational schema is provided in Appendix~\ref{sec:appendix-relational}.

\subsection{Meaningful Examples}

The northCoder repository~\cite{northcoder-repo} provides multiple \texttt{.csv} files containing a sampled dataset (including 174 titles) together with reference tables (genres, categories, languages, regions, etc.). We load these files into the SQL relational database using the provided \texttt{.sql} script and verify successful population using row-count checks per table.

% As mentioned earlier, this database will then be enriched and completed using data provided by \textit{Wikidata}. 

\subsection{SQL Implementation}

% \textcolor{red}{Docker-compose + SQL dump + mySQL?}

The SQL implementation follows the \gls{erd} described previously and the repository northCoder~\cite{northcoder-repo}. The scripts from this repository were used as the initial implementation and then adapted to better match our conceptual model.

First, several foreign key constraints were missing in the original scripts. These constraints were added to enforce the relationships defined in the \gls{erd}. During this process, an inconsistency in the dataset was discovered: some records in the \texttt{title\_principal} table referenced \texttt{talent\_id} values that were not present in the \texttt{talent} table. Since this prevented the creation of the foreign key constraint, these invalid records were removed in order to maintain referential integrity.

Second, the structure of the principal data was not fully normalized. In the original dataset, the attribute \texttt{role\_names} may contain multiple values stored as a comma-separated string. This violates the first normal form (1NF). To address this issue, the schema was normalized by splitting the original structure into two tables: \texttt{title\_principal}, which stores the participation of a talent in a title, and \texttt{principal\_role}, which stores the individual role names associated with each participation record.

These modifications ensure that the relational schema is consistent with the \gls{erd} and respects normalization and referential integrity constraints.







